% Blueprint: minimum-modulus lower bound for the ξ-canonical product
%
%   ∃ C'>0, ∀ R≥2, ∃ R'∈[R,2R], ∀ s with |s|=R':  log|P(s)| ≥ -C' R^{3/2},
%
% where P(s) = ∏_i E₁(s/ρ_i) is the genus-one canonical product over the
% multiplicity enumeration (ρ_i) of the nontrivial zeros of ζ (= zeros of ξ).
% The main theorem `thm:cm_minmod` is the statement of `lem:hf_minmod` in
% `hadamard_factorization.tex` (\Cref{cha:hadamard}).
%
% By the reverse triangle inequality `‖s-ρ‖ ≥ |‖s‖-‖ρ‖|`, a circle `|s|=R'` lies at
% distance `≥ δ` from a zero `ρ` when the real number `R'` lies at distance `≥ δ`
% from `|ρ|`. Choosing the circle therefore reduces to a one-dimensional
% Lebesgue-measure pigeonhole on `[R,2R]`.
%
% References: Levin, *Lectures on Entire Functions*, Lecture II \cite{Levin1996};
% Titchmarsh, *The theory of functions* (2nd ed.), §8.71 \cite{TitchmarshFunctions}.

\chapter{Minimum modulus of the \texorpdfstring{$\xi$}{xi}-canonical product}
\label{cha:cartan}

Throughout, $\rho\colon\iota\to\mathbb C$ is the multiplicity enumeration
\texttt{riemannXiZeros} of the nontrivial zeros of $\zeta$ (equivalently of
$\xi$), with $\iota=\texttt{RiemannXiZerosIndex}$, each $\rho_i\neq0$, and
\[
  P(s)\;=\;\texttt{riemannXiCanonicalProduct}\,\rho\,s\;=\;\prod_i E_1(s/\rho_i),
  \qquad E_1(w)=(1-w)e^{w}
\]
the genus-one canonical product of \Cref{def:hf_prod}. We write $n(t)$ for the zero
counting function \texttt{riemannXiZeroCount}\,$t$ of \Cref{def:zs_count} and recall
$n(t)\le C_n\,t\log t$ for $t\ge2$ (\Cref{lem:zs_count_bound}). All circles are centred at
$0$. The goal of the chapter is the minimum-modulus selection \Cref{thm:cm_minmod}.

\section{Elementary lower bounds for the factor and the logarithm}
\label{sec:cm_elem}

\begin{lemma}[Quadratic lower bound on the elementary factor]
\label{lem:cm_elem_lower}
\lean{neg_le_log_norm_weierstrassE1}
\leanok
\uses{lem:hf_elem_bound}
For every $w\in\mathbb C$ with $\|w\|\le\tfrac12$,
\[
  -4\,\|w\|^{2}\;\le\;\log\|E_1(w)\|.
\]
\end{lemma}

\begin{proof}
\leanok
\uses{lem:hf_elem_bound}
By \Cref{lem:hf_elem_bound}, $\|E_1(w)-1\|\le2\|w\|^{2}$, so
$\|E_1(w)\|\ge1-2\|w\|^{2}\ge1-2\cdot\tfrac14=\tfrac12>0$ (in particular $E_1(w)\neq0$).
Set $x:=2\|w\|^{2}\in[0,\tfrac12]$, so $\|E_1(w)\|\ge1-x>0$ and
$\log\|E_1(w)\|\ge\log(1-x)$. The elementary inequality $\log(1-x)\ge-2x$ for
$x\in[0,\tfrac12]$ (from $\log(1-x)\ge1-\tfrac1{1-x}=-\tfrac{x}{1-x}\ge-2x$, using
$\log y\ge1-y^{-1}$ at $y=1-x$ and $1-x\ge\tfrac12$) gives
$\log\|E_1(w)\|\ge-2x=-4\|w\|^{2}$.
\end{proof}

\begin{lemma}[A square-root domination of the squared logarithm]
\label{lem:cm_logsq_le}
\lean{Real.log_add_one_sq_le}
\leanok
For every real $R\ge1$,
\[
  (\log R+1)^{2}\;\le\;16\,R^{1/2}.
\]
Consequently $R(\log R+1)^{2}\le16\,R^{3/2}$ for all $R\ge1$.
\end{lemma}

\begin{proof}
\leanok
Using $\log(R^{1/4})=\tfrac14\log R$ and $\log y\le y-1$
(\texttt{Real.log\_le\_sub\_one\_of\_pos}) at $y=R^{1/4}>0$,
\[
  \log R=4\log(R^{1/4})\le4\bigl(R^{1/4}-1\bigr)=4R^{1/4}-4,
\]
so $\log R+1\le4R^{1/4}-3\le4R^{1/4}$. For $R\ge1$ both sides of $0\le\log R+1\le4R^{1/4}$ are
nonnegative, so squaring preserves the inequality:
$(\log R+1)^{2}\le16\,(R^{1/4})^{2}=16\,R^{1/2}$. Multiplying by $R\ge0$ gives
$R(\log R+1)^{2}\le16\,R^{3/2}$.
\end{proof}

\section{Hölder bounds on inverse zero moduli}
\label{sec:cm_sums}

The two inverse-modulus estimates of this section bound a head sum of $\|\rho_i^{-1}\|$ and a
tail sum of $\|\rho_i^{-1}\|^{2}$ by powers $t^{\pm1/2}$ of the threshold, using the
summability of the $\tfrac32$-power $S:=\sum_i\|\rho_i^{-1}\|^{3/2}$ (a special case of
the divisor-weighted summability \Cref{lem:zs_summable_divisor}) together with the elementary
monotonicities $\|\rho_i^{-1}\|=\|\rho_i^{-1}\|^{3/2}\|\rho_i\|^{1/2}$ and
$\|\rho_i^{-1}\|^{2}=\|\rho_i^{-1}\|^{3/2}\|\rho_i\|^{-1/2}$. These $t^{\pm1/2}$ rates
suffice for the order-$3/2$ bound of \Cref{thm:cm_minmod}. We first record the cardinality bound of the truncated index
set (used in the near-part estimate) and the index-type summability of the $\tfrac32$-power.

\begin{lemma}[Cardinality of the truncated zero set]
\label{lem:cm_card}
\lean{exists_card_riemannXiZeros_le}
\leanok
\uses{lem:zs_count_bound}
For each $t\ge0$ the set $I(t):=\{i\in\iota:\|\rho_i\|\le t\}$ is finite, and there is a
constant $C_n>0$ with $\#I(t)\le C_n\,t\log t$ for all $t\ge2$.
\end{lemma}

\begin{proof}
\leanok
\uses{lem:zs_count_bound}
$I(t)$ is finite because $\rho$ enumerates, with multiplicity, the zeros of $\xi$ in
$\overline B(0,t)$, of which there are finitely many of finite order (the divisor of $\xi$
has finite support on the compact $\overline B(0,t)$). Under the enumeration
$\iota=(z\in\operatorname{supp}D)\times\operatorname{Fin}((Dz).\mathtt{toNat})$ of
\Cref{lem:zs_product}, one has $\#I(t)=\sum_{\|z\|\le t}(Dz).\mathtt{toNat}
=\sum_{\|z\|\le t}(Dz)=n(t)$ by reindexing the fibre product (the fibre over $z$ has exactly
$(Dz).\mathtt{toNat}$ elements, and $(Dz).\mathtt{toNat}=Dz$ as reals since $D\ge0$ by
\texttt{divisor\_riemannXi\_nonneg}), where $D$ is the divisor of $\xi$ and
$n=\texttt{riemannXiZeroCount}$. The bound is then \Cref{lem:zs_count_bound}.
\end{proof}

\begin{lemma}[Summability of the $\tfrac32$-power of inverse moduli]
\label{lem:cm_summable_rpow}
\lean{summable_inv_norm_rpow_riemannXiZeros}
\leanok
\uses{lem:zs_summable_divisor, lem:zs_product}
For every real $p>1$ the family $i\mapsto\|\rho_i^{-1}\|^{p}$ over $\iota$ is summable.
\end{lemma}

\begin{proof}
\leanok
\uses{lem:zs_summable_divisor, lem:zs_product}
Under the fibre enumeration
$\iota=(z\in\operatorname{supp}D)\times\operatorname{Fin}((Dz).\mathtt{toNat})$ of
\Cref{lem:zs_product}, summing over each finite fibre reduces summability of
$i\mapsto\|\rho_i^{-1}\|^{p}$ to summability of $z\mapsto(Dz)\,\|z\|^{-p}$ over
$\operatorname{supp}D$, which is the divisor-weighted summability \Cref{lem:zs_summable_divisor}
for $p>1$ (here $\|\rho_i^{-1}\|^{p}=\|z\|^{-p}$ on the fibre over $z$, which has $(Dz).\mathtt{toNat}$
elements, with $(Dz).\mathtt{toNat}=Dz$ as reals since $D\ge0$). This is the proof of
\Cref{lem:zs_product} with the exponent $2$ replaced by $p$.
\end{proof}

\begin{lemma}[Head sum of inverse moduli]
\label{lem:cm_sum_inv}
\lean{exists_sum_inv_norm_riemannXiZeros_le}
\leanok
\uses{lem:cm_summable_rpow}
Writing $S:=\sum_i\|\rho_i^{-1}\|^{3/2}$ (finite by \Cref{lem:cm_summable_rpow}), there is a
constant $C_1>0$ such that for every $t\ge0$,
\[
  \sum_{i\,:\,\|\rho_i\|\le t}\|\rho_i^{-1}\|\;\le\;C_1\,t^{1/2}.
\]
\end{lemma}

\begin{proof}
\leanok
\uses{lem:cm_summable_rpow}
Take $C_1:=S+1>0$. Each $\rho_i\neq0$, so for $\|\rho_i\|\le t$ the monotonicity
$\|\rho_i\|^{1/2}\le t^{1/2}$ gives
$\|\rho_i^{-1}\|=\|\rho_i^{-1}\|^{3/2}\|\rho_i\|^{1/2}\le t^{1/2}\|\rho_i^{-1}\|^{3/2}$. Summing
over the finite head index set $\{i:\|\rho_i\|\le t\}$ and bounding the partial sum of the
nonnegative summable family $\|\rho_i^{-1}\|^{3/2}$ by its total $S$ (\Cref{lem:cm_summable_rpow},
\texttt{sum\_le\_tsum}),
\[
  \sum_{i:\|\rho_i\|\le t}\|\rho_i^{-1}\|\le t^{1/2}\sum_{i:\|\rho_i\|\le t}\|\rho_i^{-1}\|^{3/2}
   \le t^{1/2}\,S\le C_1\,t^{1/2}.
\]
\end{proof}

\begin{lemma}[Tail sum of inverse squared moduli]
\label{lem:cm_sum_inv_sq_tail}
\lean{exists_sum_inv_norm_sq_riemannXiZeros_tail_le}
\leanok
\uses{lem:cm_summable_rpow, lem:zs_product}
With $S$ as above, there is a constant $C_2>0$ such that for every $t>0$,
\[
  \sum_{i\,:\,\|\rho_i\|> t}\|\rho_i^{-1}\|^{2}\;\le\;C_2\,t^{-1/2}.
\]
\end{lemma}

\begin{proof}
\leanok
\uses{lem:cm_summable_rpow, lem:zs_product}
Take $C_2:=S+1>0$. For $\|\rho_i\|>t>0$ the monotonicity $\|\rho_i\|^{-1/2}\le t^{-1/2}$ gives
$\|\rho_i^{-1}\|^{2}=\|\rho_i^{-1}\|^{3/2}\|\rho_i\|^{-1/2}\le t^{-1/2}\|\rho_i^{-1}\|^{3/2}$. Both
$i\mapsto\|\rho_i^{-1}\|^{2}$ (\Cref{lem:zs_product}) and $i\mapsto\|\rho_i^{-1}\|^{3/2}$
(\Cref{lem:cm_summable_rpow}) are summable, so summing the termwise bound over the tail
$\{i:\|\rho_i\|>t\}$ (\texttt{tsum\_le\_tsum}) and bounding the tail of the nonnegative summable
family by its total $S$ (\texttt{Summable.tsum\_le\_tsum} over the subtype),
\[
  \sum_{i:\|\rho_i\|>t}\|\rho_i^{-1}\|^{2}\le t^{-1/2}\sum_{i:\|\rho_i\|>t}\|\rho_i^{-1}\|^{3/2}
   \le t^{-1/2}\,S\le C_2\,t^{-1/2}.
\]
\end{proof}

\section{One-dimensional radius selection}
\label{sec:cm_radius}

\begin{lemma}[Avoiding finitely many points on an interval]
\label{lem:cm_radial_select}
\lean{exists_mem_Icc_forall_le_dist}
\leanok
Let $S$ be a finite subset of $\mathbb R$ (or a finite multiset/\texttt{Finset}), let
$a\le b$ and $\delta>0$ satisfy $2\delta\,(\#S)<b-a$. Then there exists $r\in[a,b]$ with
$\delta\le|r-t|$ for every $t\in S$.
\end{lemma}

\begin{proof}
\leanok
The "bad" set $\mathcal B:=\bigcup_{t\in S}(t-\delta,t+\delta)$ has Lebesgue measure
$\le\sum_{t\in S}\operatorname{vol}(t-\delta,t+\delta)=2\delta\,(\#S)<b-a
=\operatorname{vol}[a,b]$. Hence $[a,b]\setminus\mathcal B$ has positive measure, so is
nonempty; any $r$ in it lies in $[a,b]$ and satisfies $r\notin(t-\delta,t+\delta)$, i.e.
$|r-t|\ge\delta$, for every $t\in S$.
\end{proof}

\section{Lower bounds on the near and far parts}
\label{sec:cm_parts}

We split $P(s)=\prod_iE_1(s/\rho_i)$ at the threshold $4R$: $\mathrm{near}=\{i:\|\rho_i\|\le4R\}$
and $\mathrm{far}=\{i:\|\rho_i\|>4R\}$. For $\|s\|\le2R$ the far factors have small argument and are
handled pointwise; the near factors are handled on the selected circle.

\begin{lemma}[Lower bound on the far part]
\label{lem:cm_far_lower}
\lean{exists_tsum_far_log_norm_weierstrassE1_ge}
\leanok
\uses{lem:cm_elem_lower, lem:cm_sum_inv_sq_tail}
There is a constant $C_F>0$ such that for every $R\ge2$ and every $s$ with $\|s\|\le2R$,
\[
  \sum_{i\,:\,\|\rho_i\|>4R}\log\|E_1(s/\rho_i)\|\;\ge\;-C_F\,R^{3/2}.
\]
\end{lemma}

\begin{proof}
\leanok
\uses{lem:cm_elem_lower, lem:cm_sum_inv_sq_tail}
Take $C_F:=8C_2$ with $C_2$ from \Cref{lem:cm_sum_inv_sq_tail}. For $i\in\mathrm{far}$,
$\|s/\rho_i\|=\|s\|/\|\rho_i\|\le2R/(4R)=\tfrac12$, so \Cref{lem:cm_elem_lower} gives
$\log\|E_1(s/\rho_i)\|\ge-4\|s/\rho_i\|^{2}=-4\|s\|^{2}\|\rho_i^{-1}\|^{2}
\ge-16R^{2}\|\rho_i^{-1}\|^{2}$. Both the family of logs (\Cref{lem:hf_prod_summable} via
\texttt{summable\_log\_norm\_weierstrassE1}) and the dominating family
$i\mapsto-16R^{2}\|\rho_i^{-1}\|^{2}$ (\Cref{lem:zs_product},
\texttt{summable\_inv\_norm\_sq\_riemannXiZeros}) are summable on the tail, so by
\texttt{tsum\_le\_tsum} we may sum the termwise bound:
\[
  \sum_{i\,:\,\|\rho_i\|>4R}\log\|E_1(s/\rho_i)\|
   \ge-16R^{2}\sum_{i\,:\,\|\rho_i\|>4R}\|\rho_i^{-1}\|^{2}
   \ge-16R^{2}\cdot C_2(4R)^{-1/2}=-8C_2R^{3/2},
\]
the middle inequality by \Cref{lem:cm_sum_inv_sq_tail} at threshold $4R$, and
$16R^{2}(4R)^{-1/2}=16\cdot4^{-1/2}R^{3/2}=8R^{3/2}$.
\end{proof}

\begin{lemma}[Lower bound on the near part on a good circle]
\label{lem:cm_near_lower}
\lean{exists_finset_near_log_norm_weierstrassE1_ge}
\leanok
\uses{lem:cm_card, lem:cm_sum_inv, lem:cm_radial_select, lem:cm_logsq_le}
There is a constant $C_N>0$ such that for every $R\ge2$ there is $R'\in[R,2R]$ for which the
following hold for all $s$ with $\|s\|=R'$:
\begin{enumerate}
\item $s\neq\rho_i$ for every $i\in\mathrm{near}=\{i:\|\rho_i\|\le4R\}$ (no near factor vanishes); and
\item $\displaystyle\sum_{i\,:\,\|\rho_i\|\le4R}\log\|E_1(s/\rho_i)\|\;\ge\;-C_N\,R^{3/2}.$
\end{enumerate}
\end{lemma}

\begin{proof}
\leanok
\uses{lem:cm_card, lem:cm_sum_inv, lem:cm_radial_select, lem:cm_logsq_le}
Write $m:=\#\mathrm{near}=\#\{i:\|\rho_i\|\le4R\}$ (finite by \Cref{lem:cm_card}) and
$\delta_R:=R/(4(m+1))>0$. Apply \Cref{lem:cm_radial_select} to the finite set
$S:=\{\|\rho_i\|:i\in\mathrm{near}\}$ on $[a,b]=[R,2R]$ with $\delta=\delta_R$: since
$\#S\le m$ and $2\delta_R\,m=\tfrac{R}{2}\cdot\tfrac{m}{m+1}<\tfrac R2<R=b-a$, there is
$R'\in[R,2R]$ with $|R'-\|\rho_i\||\ge\delta_R$ for every $i\in\mathrm{near}$. Fix $s$ with
$\|s\|=R'$. The reverse triangle inequality \texttt{abs\_norm\_sub\_norm\_le} gives
$\|s-\rho_i\|\ge\bigl|\|s\|-\|\rho_i\|\bigr|=|R'-\|\rho_i\||\ge\delta_R>0$ for near $i$, so
$s\neq\rho_i$, proving (1).

For (2), each near $\rho_i$ satisfies $s/\rho_i\neq1$, so by \texttt{log\_norm\_weierstrassE1},
\[
  \log\|E_1(s/\rho_i)\|=\log\|1-s/\rho_i\|+\operatorname{Re}(s/\rho_i)
   =\bigl(\log\|s-\rho_i\|-\log\|\rho_i\|\bigr)+\operatorname{Re}(s/\rho_i).
\]
Bound each summand below:
\begin{itemize}
\item $\log\|s-\rho_i\|\ge\log\delta_R$ (monotonicity of $\log$, from $\|s-\rho_i\|\ge\delta_R>0$);
\item $-\log\|\rho_i\|\ge-\log(4R)$ (from $0<\|\rho_i\|\le4R$);
\item $\operatorname{Re}(s/\rho_i)\ge-\|s/\rho_i\|=-R'\|\rho_i^{-1}\|\ge-2R\|\rho_i^{-1}\|$.
\end{itemize}
Summing over the $m$ near indices and using \Cref{lem:cm_sum_inv} at $t=4R$
($\sum_{\mathrm{near}}\|\rho_i^{-1}\|\le C_1(4R)^{1/2}$):
\[
  \sum_{\mathrm{near}}\log\|E_1(s/\rho_i)\|
   \ge m\bigl(\log\delta_R-\log(4R)\bigr)-2R\,C_1(4R)^{1/2}.
\]
The second term is $-4C_1R^{3/2}$. For the first,
$\log\delta_R-\log(4R)=\log R-\log\bigl(4(m+1)\bigr)-\log(4R)=-\log4-\log\bigl(4(m+1)\bigr)\le0$,
and $m\le C_n4R\log(4R)$ (\Cref{lem:cm_card}, as $4R\ge2$) with
$\log\bigl(4(m+1)\bigr)\le C(\log R+1)$ and $\log(4R)\le C(\log R+1)$ for $R\ge2$ gives
\[
  m\bigl(\log\delta_R-\log(4R)\bigr)\ge-m\,C(\log R+1)\ge-4C_nC\,R\log(4R)(\log R+1)
   \ge-C''(\log R+1)^{2}R.
\]
Finally $(\log R+1)^{2}R\le16R^{3/2}$ by \Cref{lem:cm_logsq_le}, so both terms are
$\ge-C_NR^{3/2}$ after enlarging the constant.
\end{proof}

\section{Assembly: the minimum-modulus selection}
\label{sec:cm_main}

\begin{theorem}[Minimum-modulus lower bound for the canonical product]
\label{thm:cm_minmod}
\lean{exists_circle_log_norm_riemannXiCanonicalProduct_ge}
\leanok
\uses{lem:cm_far_lower, lem:cm_near_lower}
There is a constant $C'>0$ such that for every $R\ge2$ there is $R'\in[R,2R]$ with
\[
  \log\|P(s)\|\;\ge\;-C'\,R^{3/2}\qquad\text{for all }s\text{ with }\|s\|=R'.
\]
This is exactly \Cref{lem:hf_minmod} of \Cref{cha:hadamard}.
\end{theorem}

\begin{proof}
\leanok
\uses{lem:cm_far_lower, lem:cm_near_lower}
Let $C_F,C_N$ be the constants of \Cref{lem:cm_far_lower} and \Cref{lem:cm_near_lower} and set
$C':=C_F+C_N>0$. Fix $R\ge2$ and let $R'\in[R,2R]$ be the radius of \Cref{lem:cm_near_lower};
note $\|s\|=R'$ implies $\|s\|\le2R$. For such $s$, part (1) of \Cref{lem:cm_near_lower} gives
$s\neq\rho_i$ for every near $i$, while every far $i$ has $\|\rho_i\|>4R\ge2R\ge\|s\|$, so
$s\neq\rho_i$ as well; thus $s\neq\rho_i$ for all $i$. Consequently $P(s)\neq0$ and, by
\texttt{log\_norm\_riemannXiCanonicalProduct}, the family $i\mapsto\log\|E_1(s/\rho_i)\|$ is
summable with
\[
  \log\|P(s)\|=\sum_i\log\|E_1(s/\rho_i)\|
   =\sum_{\|\rho_i\|\le4R}\log\|E_1(s/\rho_i)\|+\sum_{\|\rho_i\|>4R}\log\|E_1(s/\rho_i)\|,
\]
splitting the summable family (\texttt{tsum\_subtype\_add\_tsum\_subtype\_compl}) over the set
$\{i:\|\rho_i\|\le4R\}$ and its complement. Apply \Cref{lem:cm_near_lower}(2) to the first sum
and \Cref{lem:cm_far_lower} to the second:
\[
  \log\|P(s)\|\ge-C_NR^{3/2}-C_FR^{3/2}=-C'R^{3/2}\qquad(\|s\|=R'),
\]
as required.
\end{proof}

\begin{remark}
\Cref{thm:cm_minmod} is the statement of \Cref{lem:hf_minmod} of \Cref{cha:hadamard}
(\texttt{exists\_circle\_log\_norm\_riemannXiCanonicalProduct\_ge}). The classical proof uses the
two-dimensional Boutroux–Cartan exceptional-disc covering; the proof here selects the radius with
the one-dimensional \Cref{lem:cm_radial_select} and the reverse triangle inequality. The
inverse-modulus inputs \Cref{lem:cm_sum_inv} and \Cref{lem:cm_sum_inv_sq_tail} are Hölder bounds
against the summable $\tfrac32$-power series. The arithmetic inputs are the counting bound
\Cref{lem:zs_count_bound} and the zero summability \Cref{lem:zs_product} and
\Cref{lem:zs_summable_divisor} of \Cref{cha:zero_summability}; the rest is elementary real and
complex analysis from mathlib or \Cref{cha:hadamard}. The Hadamard factorisation
\Cref{thm:hf_hadamard}, and through it the lower bound \Cref{lem:lb_pf_hadamard} of
\Cref{cha:lehman_bound}, depends on this result.
\end{remark}
